\section{Reduced root data of semisimple rank \(2\)}
\label{sga3:XXI:4}

\renewcommand{\thetheorem}{\thesection.\arabic{theorem}}
\setcounter{theorem}{0}

\SGARefTarget{sga3:XXI:4.0}
\noindent\textbf{4.0.}\footnote{Editors' note: The numbering \(4.0\)
was added for later references.}
Let \(\mathcal{R}\) be a root datum of semisimple rank \(2\). Let
\(\{\alpha,\beta\}\) be a system of simple roots. Suppose that
\(\ell(\alpha)\le\ell(\beta)\). By
\ref{sga3:XXI:2.3.1} and \ref{sga3:XXI:3.2.1}, there are then four
possibilities:
\[
\begin{array}{c|c|c|c|c}
\text{Type}
 &\ell(\beta)/\ell(\alpha)
 &\ell(\beta^\ast)/\ell(\alpha^\ast)
 &(\beta^\ast,\alpha)
 &(\alpha^\ast,\beta)\\ \hline
A_1\times A_1&-&-&0&0\\
A_2&1&1&-1&-1\\
B_2&2&1/2&-1&-2\\
G_2&3&1/3&-1&-3
\end{array}
\]
It follows from Remark~\ref{sga3:XXI:3.4.5} that the order of
\(s_\alpha s_\beta\) is respectively \(2,3,4,6\).

We study each of these systems separately and give the list of
indivisible roots.

\medskip
\noindent\emph{Type \(A_1\times A_1\).}
The indivisible roots are
\[
 \alpha,\quad\beta,\quad-\alpha,\quad-\beta.
\]
The corresponding coroots are
\[
 \alpha^\ast,\quad\beta^\ast,\quad-\alpha^\ast,\quad-\beta^\ast.
\]

\medskip
\noindent\emph{Type \(A_2\).}
The indivisible positive roots are as follows:
\[
\begin{array}{c|ccc}
\text{root }\gamma
 &\alpha&\beta&\alpha+\beta\\ \hline
\ell(\gamma)/\ell(\alpha)
 &1&1&1\\ \hline
\text{coroot }\gamma^\ast
 &\alpha^\ast&\beta^\ast&\alpha^\ast+\beta^\ast\\ \hline
\ell(\gamma^\ast)/\ell(\beta^\ast)
 &1&1&1
\end{array}
\]
The half-sum of the indivisible positive roots is
\[
 \rho=\alpha+\beta.
\]

\medskip
\noindent\emph{Type \(B_2\).}
The indivisible positive roots are as follows:
\[
\begin{array}{c|cccc}
\text{root }\gamma
 &\alpha&\beta&\alpha+\beta&2\alpha+\beta\\ \hline
\ell(\gamma)/\ell(\alpha)
 &1&2&1&2\\ \hline
\text{coroot }\gamma^\ast
 &\alpha^\ast&\beta^\ast
 &2\alpha^\ast+\beta^\ast&\alpha^\ast+\beta^\ast\\ \hline
\ell(\gamma^\ast)/\ell(\beta^\ast)
 &2&1&2&1
\end{array}
\]
The half-sum of the indivisible positive roots is
\[
 \rho=\frac{4\alpha+3\beta}{2}.
\]

\medskip
\noindent\emph{Type \(G_2\).}
The indivisible positive roots are as follows:
\begin{center}
\resizebox{\textwidth}{!}{%
$\displaystyle
\begin{array}{c|cccccc}
\text{root }\gamma
 &\alpha&\beta&\alpha+\beta&2\alpha+\beta
 &3\alpha+\beta&3\alpha+2\beta\\ \hline
\ell(\gamma)/\ell(\alpha)
 &1&3&1&1&3&3\\ \hline
\text{coroot }\gamma^\ast
 &\alpha^\ast&\beta^\ast&\alpha^\ast+3\beta^\ast
 &2\alpha^\ast+3\beta^\ast&\alpha^\ast+\beta^\ast
 &\alpha^\ast+2\beta^\ast\\ \hline
\ell(\gamma^\ast)/\ell(\beta^\ast)
 &3&1&3&3&1&1
\end{array}
$%
}
\end{center}
The half-sum of the indivisible positive roots is
\[
 \rho=5\alpha+3\beta.
\]
